\chapter*{الملخص}
\addcontentsline{toc}{chapter}{الملخص}

\selectlanguage{arabic}

هذا التقرير يعرض ملخص العمل الذي تم إنجازه في إطار مشروع نهاية الدراسة الموسوم ب" تحديث أداة إدارة متكاملة لخدمات المحاسبة" بشركة Tamtam International، لاستكمال متطلبات الحصول على شهادة الهندسة في أنظمة الشبكات والمعلومات بكلية العلوم والتقنيات بمراكش، جامعة القاضي عياض.

الهدف من هذا المشروع هو تحديث أداة إدارة متكاملة مصممة لمراقبة خدمات المحاسبة، وتكييفها مع النظام البيئي الموجه للمهن الاقتصادية والقانونية.

ركز عملي على قسمين رئيسيين: من جهة، تسهيل استشارة الخدمات المعالجة، ومن جهة أخرى، تحضير الفوترة الشهرية من خلال تطبيق آلية تحييد. يشمل ذلك توليد التقارير الشهرية للخدمات، وتكوين مستويات كبار الموظفين، وإعداد أكواد OAM.

لإدارة المشروع، اعتمدنا منهجية Agile SCRUM، حيث قسمنا العمل إلى sprints. يعتمد التنفيذ على تقنيات ويب حديثة لضمان جودة المشروع، وهي PHP مع إطار عمل Symfony للواجهة الخلفية، و React.js للواجهة الأمامية، و MySQL كنظام إدارة قواعد البيانات.

\textbf{الكلمات المفتاحية:} SCRUM، التكامل، الاستيراد والمزامنة، العمارة السادسة الأضلاع

\textbf{التقنيات:} PHP، Symfony، ReactJS، Redux، MySQL

\vspace{0.5cm}

\begin{flushright}
    \textit{EL MOUKTADIR Mohamed} \\
    \textit{شهادة الهندسة - تخصص تكنولوجيا المعلومات} \\
    \textit{السنة الجامعية: 2025/2026}
end{flushright}

\selectlanguage{french}
